%%%%%%%%%%%%%%%%
%%% gen 16 %%%
%%%%%%%%%%%%%%%%

\Chapter{16}

\Verse{1}
{
	\hJ{וְשָׂרַי אֵשֶׁת אַבְרָם לֹא יָלְדָה לוֹ וְלָהּ שִׁפְחָה מִצְרִית וּשְׁמָהּ הָגָר׃}
}
{
	\eJ{
		And Abram’s wife Sarai had not given birth by him.
		And she had an Egyptian maid, and her name was Hagar.
	}
}
{
	Previously on, bible:
	
	YHWH promised Abe infinite {\pussyb}.
	
	This time on, bible:
	
	Abe's wife wasn't even pregnant yet.
	
	They had a maid from Egypt though.
}

\Verse{2}
{
	\hJ{וַתֹּאמֶר שָׂרַי אֶל אַבְרָם הִנֵּה נָא עֲצָרַנִי יְהֹוָה מִלֶּדֶת בֹּא נָא אֶל שִׁפְחָתִי אוּלַי אִבָּנֶה מִמֶּנָּה וַיִּשְׁמַע אַבְרָם לְקוֹל שָׂרָי׃}
}
{
	\eJ{
		And Sarai said to Abram, “Here, YHWH has held me back from
		giving birth. Come to my maid. Maybe I’ll get ‘childed’
		through her.” And Abram listened to Sarai’s voice.
	}
}
{
	Abe's wife suggests he get to know the maid.
	
	He listens.
}

\Verse{3}
{
	\hP{וַתִּקַּח שָׂרַי אֵשֶׁת אַבְרָם אֶת הָגָר הַמִּצְרִית שִׁפְחָתָהּ מִקֵּץ עֶשֶׂר שָׁנִים לְשֶׁבֶת אַבְרָם בְּאֶרֶץ כְּנָעַן וַתִּתֵּן אֹתָהּ לְאַבְרָם אִישָׁהּ לוֹ לְאִשָּׁה׃}
}
{
	\eP{
		And Sarai, Abram’s wife, took Hagar, the Egyptian, her
		maid, at the end of ten years of Abram’s living in the
		land of Canaan, and gave her to Abram, her husband, as a
		wife to him.
	}
}
{
	So Abe's wife gives her maid to Abe, her husband, as a wife to him.
}

\Verse{4}
{
	\hJ{וַיָּבֹא אֶל הָגָר וַתַּהַר וַתֵּרֶא כִּי הָרָתָה וַתֵּקַל גְּבִרְתָּהּ בְּעֵינֶיהָ׃}
}
{
	\eJ{
		And he came to Hagar, and she became pregnant. And she saw
		that she had become pregnant, and her mistress was lowered
		in her eyes.
	}
}
{
	Abe comes \redacted{to} her, and she's pregnant.
	
	And she saw that she was pregnant.
	
	New wife looks down on old wife, which is mean.
}

\Verse{5}
{
	\hJ{וַתֹּאמֶר שָׂרַי אֶל אַבְרָם חֲמָסִי עָלֶיךָ אָנֹכִי נָתַתִּי שִׁפְחָתִי בְּחֵיקֶךָ וַתֵּרֶא כִּי הָרָתָה וָאֵקַל בְּעֵינֶיהָ יִשְׁפֹּט יְהֹוָה בֵּינִי וּבֵינֶיךָ׃}
}
{
	\eJ{
		And Sarai said to Abram, “My injury is on you. I, I,
		placed my maid in your bosom, and she saw that she had
		become pregnant, and I was lowered in her eyes. Let YHWH
		judge between me and you.”
	}
}
{
	Sarai blames Abe.
}

\Verse{6}
{
	\hJ{וַיֹּאמֶר אַבְרָם אֶל שָׂרַי הִנֵּה שִׁפְחָתֵךְ בְּיָדֵךְ עֲשִׂי לָהּ הַטּוֹב בְּעֵינָיִךְ וַתְּעַנֶּהָ שָׂרַי וַתִּבְרַח מִפָּנֶיהָ׃}
}
{
	\eJ{
		And Abram said to Sarai, “Here, your maid is in your hand.
		Do to her whatever is good in your eyes.” And Sarai
		degraded her, and she fled from her.
	}
}
{
So Abe's like ``Look babe I defer to you. This was your idea, and by the way it was a really great one.

I hope you keep having more ideas like this.

No complaints from me, do whatever you want.''

So Sarai degrades new wife.

It was her idea but she's jealous now.

Now Hagar's like ``Screw you guys I'm gonna go live in the woods.''
%	\footnote{
%		And Sarai degraded her, and she fled. Sarai’s treatment of the Egyptian
%		Hagar foreshadows (or is reversed in, or governs) Israel’s experience in
%		Egypt. These exact words recur there: Egypt degraded them (Exod 1:12),
%		and they fled (14:5). And note that, like Israel, Hagar fled to the
%		wilderness (v. 7)
%	}
}

\Verse{7}
{
	\hJ{וַיִּמְצָאָהּ מַלְאַךְ יְהֹוָה עַל עֵין הַמַּיִם בַּמִּדְבָּר עַל הָעַיִן בְּדֶרֶךְ שׁוּר׃}
}
{
	\eJ{
		And an angel of YHWH found her at a spring of water
		in the wilderness, by the spring on the way to Shur,
	}
}
{
	So YHWH sends a messenger 
}

\Verse{8}
{
	\hJ{וַיֹּאמַר הָגָר שִׁפְחַת שָׂרַי אֵי מִזֶּה בָאת וְאָנָה תֵלֵכִי וַתֹּאמֶר מִפְּנֵי שָׂרַי גְּבִרְתִּי אָנֹכִי בֹּרַחַת׃}
}
{
	\eJ{
		and said, “Hagar, Sarai’s maid, from where have you come,
		and where will you go?” And she said, “I’m fleeing from
		Sarai, my mistress.”
	}
}
{
}

\Verse{9}
{
	\hJ{וַיֹּאמֶר לָהּ מַלְאַךְ יְהֹוָה שׁוּבִי אֶל גְּבִרְתֵּךְ וְהִתְעַנִּי תַּחַת יָדֶיהָ׃}
}
{
	\eJ{
		And the angel of YHWH said to her, “Go back to your
		mistress, and suffer the degradation under her hands.”
	}
}
{
}

\Verse{10}
{
	\hJ{וַיֹּאמֶר לָהּ מַלְאַךְ יְהֹוָה הַרְבָּה אַרְבֶּה אֶת זַרְעֵךְ וְלֹא יִסָּפֵר מֵרֹב׃}
}
{
	\eJ{
		And the angel of YHWH said, “I’ll multiply your seed, and
		it won’t be countable because of its great number.”
	}
}
{
}

\Verse{11}
{
	\hJ{וַיֹּאמֶר לָהּ מַלְאַךְ יְהֹוָה הִנָּךְ הָרָה וְיֹלַדְתְּ בֵּן וְקָרָאת שְׁמוֹ יִשְׁמָעֵאל כִּי שָׁמַע יְהֹוָה אֶל ענְיֵךְ׃}
}
{
	\eJ{
		And the angel of YHWH said to her, “Here, you’re pregnant
		and will give birth to a son, and you shall call his name
		Ishmael, for YHWH has listened to your suffering.
	}
}
{
}

\Verse{12}
{
	\hJ{וְהוּא יִהְיֶה פֶּרֶא אָדָם יָדוֹ בַכֹּל וְיַד כֹּל בּוֹ וְעַל פְּנֵי כל אֶחָיו יִשְׁכֹּן׃}
}
{
	\eJ{
		And he’ll be a wild ass of a man, his hand against
		everyone, and everyone’s hand against him, and he’ll tent
		among all his brothers.”
	}
}
{
%	\footnote{
%		he’ll tent. The verb is denominative; that is, the verb “to tent”
%		(Hebrew root škn) derives from the word miškn, a tent or tabernacle. In
%		the blessing of Noah’s son Yaphet, it refers explicitly to dwelling in
%		tents (Gen 9:27). In Moses’ blessing of Benjamin (Deut 33:12) it is
%		connected to God’s “sheltering” (Hebrew pp, related to upph, a canopy).
%		Most important, it will be used (in its Piel form) for the way in which
%		God tents the Ark of the Testimony inside the Tabernacle. Occasionally
%		it can refer to “dwelling” in general, not necessarily in tents (e.g.,
%		Gen 49:13).
%	}
}

\Verse{13}
{
	\hJ{וַתִּקְרָא שֵׁם יְהֹוָה הַדֹּבֵר אֵלֶיהָ אַתָּה אֵל רֳאִי כִּי אָמְרָה הֲגַם הֲלֹם רָאִיתִי אַחֲרֵי רֹאִי׃}
}
{
	\eJ{
		And she called the name of YHWH who spoke to her “You are
		El-roi,” for she said, “Have I also seen after the one who
		sees me here?”
	}
}
{
}

\Verse{14}
{
	\hJ{עַל כֵּן קָרָא לַבְּאֵר בְּאֵר לַחַי רֹאִי הִנֵּה בֵין קָדֵשׁ וּבֵין בָּרֶד׃}
}
{
	\eJ{
		On account of this the well was called “the well
		Lahai-roi.” Here it is between Kadesh and Bered.
	}
}
{
}

\Verse{15}
{
	\hP{וַתֵּלֶד הָגָר לְאַבְרָם בֵּן וַיִּקְרָא אַבְרָם שֶׁם בְּנוֹ אֲשֶׁר יָלְדָה הָגָר יִשְׁמָעֵאל׃}
}
{
	\eP{
		And Hagar gave birth to a son for Abram, and Abram called
		the name of his son whom Hagar had borne Ishmael.
	}
}
{
}

\Verse{16}
{
	\hP{וְאַבְרָם בֶּן שְׁמֹנִים שָׁנָה וְשֵׁשׁ שָׁנִים בְּלֶדֶת הָגָר אֶת יִשְׁמָעֵאל לְאַבְרָם׃}
}
{
	\eP{
		And Abram was eighty years and six years old when Hagar
		gave birth to Ishmael for Abram.
	}
}
{
}
